% Ch9 WS2 -- FRQ practice on CED content, plus optional MO enrichment.
\documentclass{apchem-worksheet}

\apcunitword{Chapter}
\apcunit{9}
\apcunittitle{Covalent Bonding: Orbitals}
\apcdoctitle{Worksheet 2 \textbullet\ FRQ Practice $+$ Optional Enrichment}
\apcsubtitle{Zumdahl \S9.1 examinable \textbullet\ \S9.2--9.3 optional}

\begin{document}
\wsheader[Problems 1--2 are exam-relevant. Problem 3 is marked OPTIONAL ---
molecular orbital theory is not on the AP exam, and skipping it costs you
nothing.]

\problem[]{\textbf{FRQ (10 points).} Consider the molecule
\ce{CH3CH=CH2} (propene) and the ion \ce{NO2-}.}
\begin{parts}
  \item Give the hybridization of each of the three carbon atoms in
        propene, with the domain count that justifies each.
        \answerlines[3]{\ce{CH3} carbon: 4 domains (four single bonds) $\to$
        sp$^3$. The two doubly bonded carbons: 3 domains each $\to$ sp$^2$.}
  \item State the total number of $\sigma$ and $\pi$ bonds in propene.
        \answerlines[2]{Propene, \ce{C3H6}: 8 $\sigma$ and 1 $\pi$.}
  \item Explain why rotation is free about the \ce{CH3-CH} bond but
        restricted about the \ce{CH=CH2} bond.
        \answerlines[3]{The single bond is purely $\sigma$ and
        cylindrically symmetric, so rotation changes no overlap. The double
        bond includes a $\pi$ component whose parallel $p$ orbitals must
        stay aligned; rotating would destroy that overlap, which requires
        breaking the bond.}
  \item For \ce{NO2-} (18 valence electrons), give the electron-domain
        count, the molecular geometry, and the hybridization of nitrogen.
        \answerlines[2]{3 domains including one lone pair on N $\to$ bent
        molecular geometry, nitrogen is sp$^2$.}
\end{parts}

\begin{apcnote}[Rubric --- score yourself, 10 points]
\textbf{(a) 3 pts} 1 per carbon, each requiring the domain count as
justification --- bare labels earn half \textbullet\ \textbf{(b) 2 pts}
1 $\sigma$ count, 1 $\pi$ count \textbullet\ \textbf{(c) 3 pts} 1 names
$\sigma$ symmetry, 1 names the $\pi$ requirement for parallel $p$ orbitals,
1 concludes that rotating breaks the bond \textbullet\ \textbf{(d) 2 pts}
1 domain count with the lone pair identified, 1 geometry AND sp$^2$.
\end{apcnote}

\problem[]{Assign hybridization and predict the bond angle:}
\begin{parts}
  \item P in \ce{PH3}: \answerblank[5cm]{sp$^3$, $\approx107^\circ$}
  \item C in \ce{HCN}: \answerblank[5cm]{sp, $180^\circ$}
  \item S in \ce{SO3}: \answerblank[5cm]{sp$^2$, $120^\circ$}
  \item O in \ce{CH3OH}: \answerblank[5cm]{sp$^3$, $\approx104.5^\circ$}
\end{parts}

\problem[]{\textbf{OPTIONAL --- enrichment only.} Molecular orbital theory
is not assessed on the AP exam. Attempt this only if you worked Block 3.}
\begin{parts}
  \item Give the formula for bond order in MO theory.
        \answerblank[6cm]{(bonding $-$ antibonding) $\div$ 2}
  \item \ce{He2} would have 2 bonding and 2 antibonding electrons.
        Compute the bond order and explain what it means physically.
        \answerlines[2]{Bond order 0 --- no net stabilization, so the
        molecule does not form. This is why helium is monatomic.}
  \item The Lewis structure of \ce{O2} shows all electrons paired, yet
        liquid oxygen is attracted to a magnet. Explain the discrepancy and
        how MO theory resolves it. \answerlines[3]{Paired electrons would
        make \ce{O2} diamagnetic, but it is observed to be paramagnetic
        with two unpaired electrons --- the Lewis model is simply wrong
        here. MO theory places the last two electrons in separate
        degenerate $\pi^*$ orbitals with parallel spins, predicting bond
        order 2 \emph{and} two unpaired electrons, matching both the double
        bond and the magnetism.}
  \item State one thing each model does better than the other.
        \answerlines[2]{Localized electron model: predicts geometry and is
        far simpler to apply. MO theory: correctly predicts magnetic
        behaviour and handles delocalization.}
\end{parts}

\begin{spiralstrip}{Chapter 9 \textbullet\ block 1--2}
  \item Hybridization of C in \ce{CO2}: \answerblank[1.8cm]{sp}
  \item $\sigma$ and $\pi$ in \ce{C2H4}:
        \answerblank[3.4cm]{5 $\sigma$, 1 $\pi$}
  \item Domains on N in \ce{NH3}: \answerblank[1.6cm]{4}
  \item Which bond type blocks rotation? \answerblank[1.8cm]{$\pi$}
  \item Number of hybrid orbitals sp$^2$ produces:
        \answerblank[1.4cm]{3}
\end{spiralstrip}

\end{document}
